% !TeX program = lualatex
\documentclass[11pt,a4paper]{article}
\usepackage{geometry}
\usepackage{hyperref}
\usepackage{booktabs}
\usepackage{longtable}
\usepackage{enumitem}
\usepackage{xurl}
\usepackage{xcolor}
\usepackage{url}
\usepackage{listings}
\usepackage{fontspec}
\setmonofont{Ubuntu Mono}

\hypersetup{colorlinks=true,linkcolor=black,urlcolor=blue,pdftitle={Memoria del proyecto S\&P 500}}

\title{\LARGE Análisis y evolución de las empresas del S\&P 500\\\large Memoria del proyecto de Shell y Entornos Linux}
\author{Yassin Pellicer Lamla}
\date{\today}

\begin{document}
\maketitle
\tableofcontents
\newpage

\section{Objetivo}
El proyecto automatiza la descarga, depuración, análisis y presentación de información financiera de las compañías del índice S\&P 500. Para cada ejecución diaria se conserva una instantánea del conjunto de datos y se generan informes de texto y HTML. 

Los informes permiten consultar el ranking de capitalización bursátil, la agrupación por sectores y su ranking, las empresas que más próximas están a su máximo de 52 semanas y la evolución histórica de precio y capitalización.

El objetivo final es disponer de un sistema reproducible y automatizado que permita obtener periódicamente una visión general del estado de las empresas del S\&P 500, al mismo tiempo que se aplican herramientas y conceptos propios de Shell y de la administración de sistemas Linux, como la ejecución programada de tareas, el procesamiento de archivos y la generación automática de informes.

\section{Elección del dataset}
El dataset utilizado en este proyecto consiste en un archivo .csv de 500 entradas correspondientes a los datos de cada una de las 500 empresas más rentables de la bolsa estadounidense. En este caso se utiliza un único dataset en formato CSV publicado por el repositorio abierto \emph{S\&P 500 Companies Financials} de la organización \texttt{datasets} en GitHub. Se descarga desde:

\begin{flushleft}
\url{https://raw.githubusercontent.com/datasets/s-and-p-500-companies-financials/main/data/constituents-financials.csv}
\newline\newline
El repositorio público del dataset es:

\url{https://github.com/datasets/s-and-p-500-companies-financials}
\end{flushleft}

La fuente es pública y accesible por descarga mediante el comando \texttt{wget}. El proyecto no presupone que el archivo sea una fuente de datos en tiempo real, ya que se actualiza una vez al día: simplemente satisface el requisito de actualización mediante una descarga programada diariamente, que guarda una instantánea con fecha propia.

\subsection{Formato y estructura}
El formato es CSV con una fila de cabecera y una fila por compañía. El script de descarga lo almacena como \texttt{datasets/constituents-financials\_AAAAMMDD.csv}. Los campos son los siguientes:

{
\renewcommand{\arraystretch}{0.85}

\begin{longtable}{@{}p{0.37\textwidth}p{0.59\textwidth}@{}}
\toprule
Campo & Descripción \\
\midrule
\endhead

\texttt{Symbol} & Símbolo bursátil de la compañía. \\
\texttt{Name} & Nombre de la compañía. \\
\texttt{Sector} & Sector o industria asignada. \\
\texttt{Price} & Precio de la acción. \\
\texttt{Price/Earnings} & Relación precio-beneficio. \\
\texttt{Dividend Yield} & Rentabilidad por dividendo. \\
\texttt{Earnings/Share} & Beneficio por acción. \\
\texttt{52 Week Low} & Precio mínimo durante las últimas 52 semanas. \\
\texttt{52 Week High} & Precio máximo durante las últimas 52 semanas. \\
\texttt{Market Cap} & Capitalización bursátil. \\
\texttt{EBITDA} & Beneficio antes de intereses, impuestos, depreciaciones y amortizaciones. \\
\texttt{Price/Sales} & Relación entre el precio y las ventas. \\
\texttt{Price/Book} & Relación entre el precio de mercado y el valor contable. \\
\texttt{SEC Filings} & URL de documentación presentada ante la SEC. \\

\bottomrule
\end{longtable}
}

\subsection{Datos relevantes y preprocesamiento}

No todos estos datos son relevantes para el estudio. En el proyecto sólo se consideran los campos \texttt{Symbol, Name, Sector, Price, 52 Week High y Market Cap}. El preprocesador invocado en el servicio de descarga automático se deshace de cualquier entrada que no tenga, como mínimo, estas columnas. Sin embargo, las columnas no significativas se conservan de todas maneras para permitir la ampliación y creación de métricas más complejas en un futuro.

\section{Automatización mediante cron}
El archivo \texttt{cron/crontab.txt} define una ejecución diaria a las 12:00 horas
\begin{verbatim}
0 12 * * * /ruta/al/proyecto/cron/script_cron.sh
\end{verbatim}
El script \texttt{cron/\allowbreak script\_cron.sh} ejecuta las etapas de limpieza, descarga, saneamiento y creación de informes. El procesamiento está separado en scripts Bash con responsabilidades concretas para cada uno:
\begin{itemize}[noitemsep]
  \item \texttt{script\_delete.sh}: elimina snapshots CSV con más de 30 días.
  \item \texttt{script\_download.sh}: descarga el CSV y reintenta la operación hasta diez veces ante un fallo de red.
  \item \texttt{script\_sanitize.sh}: valida los campos requeridos y normaliza el CSV.
  \item \texttt{script\_analisis\_csv\_txt.sh}: genera los informes de texto.
  \item \texttt{script\_analisis\_csv\_html.sh}: genera el informe HTML a partir de las plantillas localizadas en \texttt{scripts/html/}.
\end{itemize}
Los scripts de análisis aceptan opcionalmente una fecha \texttt{AAAAMMDD}; si se omite, emplean la fecha actual. Esto permite regenerar un informe histórico sin volver a descargar el dataset.

\newpage
\section{Scripts de procesamiento}

\texttt{cron/\allowbreak script\_cron.sh} inicia el proceso automatizado por el servicio \texttt{cron} y ejecuta los 3 siguientes scripts de procesamiento

\begin{itemize}[noitemsep]
    \item \texttt{script\_delete.sh}
    \item \texttt{script\_download.sh}
    \item \texttt{script\_sanitize.sh}
\end{itemize}

\subsection{\texttt{script\_delete.sh}}

Este script borra por defecto los datasets con una antigüedad mayor de 30 días. Concretamente, extrae la fecha del nombre de los archivos y, mediante un bucle, comprueba si el archivo consultado tiene una antigüedad superior a 30 días respecto a la fecha actual. En ese caso, elimina el archivo.

La conservación del histórico de archivos durante 30 días permite comparar e interpretar cambios y comprender mejor la evolución de la economía durante un período corto y significativo sin generar reportes excesivamente largos.

% \lstset{
%     basicstyle=\ttfamily\small,
%     keywordstyle=\color{blue},
%     commentstyle=\color{gray},
%     stringstyle=\color{purple},
%     breaklines=true,
% }

% \begin{lstlisting}
% # ...
%   if [[ "$dataset_date" < "$(date -d "$DAYS days ago" +%Y%m%d)" ]]; then
%    rm -- "$dataset" 2> >(log_command_error "rm") && files_to_delete=$((files_to_delete + 1))
% # ...
% \end{lstlisting}

\subsection{\texttt{script\_download.sh}}

Este script descarga el .csv del remoto y le da un nombre especial que le añade la coletilla \texttt{\_<date>} al final, de manera que el .csv descargado el 11/09/2026 se llamaría \texttt{constituents-financials\_20260911.csv}. 

El script es a prueba de fallos de conexión desde el cliente. Es decir, si por algún motivo el equipo pierde la red, el programa reintentará la descarga hasta 10 veces en un plazo de 6 minutos. También controla las descargas vacías o la ausencia del recurso a descargar. 

\subsection{\texttt{script\_sanitize.sh}}

Este script toma el dataset recién descargado y comprueba que su estructura es válida para la extracción y análisis de datos posterior.

En primer lugar elimina todas aquellas entradas del .csv que carezcan de las columnas \texttt{Symbol, Name, Sector, Price, 52 Week High y Market Cap}. A continuación reformatea los campos \texttt{Name y Sector} para eliminar las ',' embebidas entre comillas '"'  que dificultan la lectura del archivo (ya que el archivo utiliza la ',' como símbolo ortográfico en nombres como "Nike, Inc.", lo que impide la correcta separación de los campos entre comas).

\section{Extracción y análisis de datos}

Una vez ejecutados satisfactoriamente los tres primeros scripts de procesamiento, se ejecutan los dos scripts correspondientes a la extracción y análisis de los datos del dataset:

\begin{itemize}[noitemsep]
\item \texttt{script\_analisis\_csv\_txt.sh}
\item \texttt{script\_analisis\_csv\_html.sh}
\end{itemize}

Ambos scripts realizan los mismos análisis sobre el dataset (y con los mismos algoritmos), pero generan los resultados en formatos diferentes.

El script \texttt{script\_analisis\_csv\_txt.sh} genera cuatro archivos de texto, cuyos nombres siguen el formato indicado a continuación:

\begin{itemize}[noitemsep]
\item \texttt{highest\_grossing\_sectors\_<date>.txt}
\item \texttt{highest\_winners\_52\_week\_metric\_<date>.txt}
\item \texttt{price\_marketcap\_evolution\_<date>.txt}
\item \texttt{top\_market\_cap\_companies\_<date>.txt}
\end{itemize}

En todos los casos, \texttt{<date>} corresponde a la fecha del día en que se generan los archivos.

Por otro lado, el script \texttt{script\_analisis\_csv\_html.sh} genera un informe en formato HTML compuesto por los siguientes archivos:

\begin{itemize}[noitemsep]
\item \texttt{report\_<date>.html}
\item \texttt{script.js}
\item \texttt{style.css}
\end{itemize}

Los reportes \texttt{.txt} se almacenan en la ruta \texttt{project/analysis/<date>} donde \texttt{<date>} es la fecha del reporte en cuestión. El HTML generado se anida en una carpeta adicional llamada \texttt{report}. Esta es la estructura de los archivos de reporte.
\begin{verbatim}
project/
└── analysis/
    └── <date>/
        ├── highest_grossing_sectors_<date>.txt
        ├── highest_winners_52_week_metric_<date>.txt
        ├── price_marketcap_evolution_<date>.txt
        ├── top_market_cap_companies_<date>.txt
        └── report/
            ├── report_<date>.html
            ├── script.js
            └── style.css
\end{verbatim}

\subsection{Reporte de empresas con mayor capitalización bursátil}
Este reporte se genera ordenando la columna de capitalización bursátil (\texttt{Market Cap}) en orden descendiente por medio del siguiente flujo:

\begin{enumerate}[noitemsep]
\item Se lee el fichero excluyendo la cabecera mediante el comando \texttt{tail}.
\item \texttt{sort} ordena las entradas por la columna \texttt{Market Cap} en orden descendente.
\end{enumerate}

\subsection{Reporte de sectores con mayor capitalización bursátil}

Este reporte agrupa las empresas según su sector y calcula la capitalización
bursátil total de cada uno mediante el siguiente flujo:

\begin{enumerate}[noitemsep]
\item Se lee el fichero sin la cabecera mediante \texttt{tail}.
\item \texttt{sort} ordena las entradas por la columna \texttt{Market Cap} en orden descendente.
\item \texttt{awk} agrupa las empresas por sector y suma sus valores de \texttt{Market Cap}.
\item Los sectores se ordenan de forma descendente según su capitalización total mediante un algoritmo de ordenación en \texttt{awk} 
\end{enumerate}

\subsection{Reporte de mayores ganadores por métrica de 52 semanas}

Este reporte identifica las empresas cuyo precio está más próximo al máximo de las últimas 52 semanas:

\begin{enumerate}[noitemsep]
\item Se lee el fichero sin la cabecera mediante \texttt{tail}.
\item \texttt{awk} calcula la distancia porcentual entre \texttt{Price} y
\texttt{52 Week High}.
\item \texttt{sort} ordena las empresas de menor a mayor distancia.
\item \texttt{cut} elimina el campo auxiliares donde se calculó esta distancia.
\end{enumerate}

\subsection{Reporte de evolución del mercado}

Este reporte muestra la evolución histórica del precio y de la capitalización bursátil de las empresas:

\begin{enumerate}[noitemsep]
\item Se recorren los datasets disponibles hasta la fecha seleccionada y mediante \texttt{awk} se vuelca cada entrada de cada empresa por día en la salida. \texttt{sort} las ordena previamente por capitlización, aunque no es necesario.
\item la salida anterior se le pasa a \texttt{awk}, que calcula la capitalización media de la empresa durante el período analizado y se la añade a cada fila suministrada.
\item \texttt{sort} ordena la entrada por prioridad, primero por \texttt{Symbol}, luego por \texttt{Name} y por último
\end{enumerate}

\section{Registro (log)}
El archivo \texttt{log.txt} conserva un registro cronológico de todas las ejecuciones. Cada línea incluye fecha y hora con milisegundos, además del estado \texttt{[ok]} o \texttt{[ko]}. Se registran el inicio y el final de cada etapa, los archivos de entrada y salida, los reintentos de descarga y los registros descartados durante el saneamiento. Este registro permite auditar la automatización y diagnosticar errores.

\section{Estructura de ficheros del proyecto}
\begin{verbatim}
project/
├── analysis/
│   ├── highest_grossing_sectors_<date>.txt
│   ├── highest_winners_52_week_metric_<date>.txt
│   ├── price_marketcap_evolution_<date>.txt
│   ├── top_market_cap_companies_<date>.txt
│   └── report/
│       ├── report_<date>.html
│       ├── script.js
│       └── style.css
├── cron/
\end{verbatim}

\newpage

\appendix
\section{Uso de IA en funciones de ayuda y formateo}
Se utilizó asistencia de IA como apoyo al desarrollo, centrada en tareas de bajo riesgo y revisables: mejora del formato del informe HTML, definición de estilos CSS, incorporación de numeración y totales en tablas, y redacción inicial de esta memoria. La IA no sustituye la validación del proyecto: las decisiones sobre el dataset, la automatización, los comandos ejecutados y la verificación de los resultados se revisaron sobre los archivos y scripts del repositorio. No se empleó para inventar datos financieros ni para modificar la fuente de datos descargada.

\end{document}
